\chapter{Происхождение центральной щели синглета и триплета из родительского действия}
\label{ch:v8-baryon-c0-singlet-triplet-central-gap-parent-action-origin-gate}

Предыдущий гейт доказал, что минимальное равновесное данное есть
$\theta=\beta\Delta$. Теперь проверим, содержится ли направление и
коэффициент $\Delta$ в уже объявленных grading-, семейных и концевых
операторах.

\section{Единое центральное направление}

На представлении $\mathbf1\oplus\mathbf3$ введём следово-центрированный
оператор
\begin{equation}
 Q=P_{\mathbf3}-\frac34I_4,
 \qquad \Tr Q=0,
 \qquad \Spec Q=\left\{-\frac34,\left(\frac14\right)^{\times3}\right\}.
 \label{eq:v8-baryon-c0-central-gap-traceless-direction}
\end{equation}
Для стандартных вещественных генераторов семейной тройки полный
квадратичный оператор Казимира равен
\begin{equation}
 \mathcal C_2=-\sum_{a=1}^{3}J_a^2=2P_{\mathbf3}.
 \label{eq:v8-baryon-c0-family-casimir-central-projector}
\end{equation}
Для минимального endpoint-расширения положим
$\Gamma_4=\operatorname{diag}(-1,+1,+1,+1)$. После центрирования получаем
точное совпадение
\begin{equation}
 \Gamma_4-\frac{\Tr\Gamma_4}{4}I_4
 =\mathcal C_2-\frac{\Tr\mathcal C_2}{4}I_4
 =2Q.
 \label{eq:v8-baryon-c0-centered-grading-casimir-equivalence}
\end{equation}
Поэтому grading, семейный оператор Казимира и центральный endpoint-проектор
порождают одну и ту же двумерную алгебру
\begin{equation}
 \operatorname{span}\{I_4,\Gamma_4,\mathcal C_2,P_{\mathbf3}\}
 =\operatorname{span}\{I_4,Q\},
 \qquad \dim=2.
 \label{eq:v8-baryon-c0-central-gap-seed-algebra}
\end{equation}
Форма центральной щели тем самым условно канонична:
\begin{equation}
 h_{\rm cen}=\varepsilon I_4+\lambda Q,
 \qquad
 E_{\mathbf3}-E_{\mathbf1}=\lambda.
 \label{eq:v8-baryon-c0-central-gap-general-hamiltonian}
\end{equation}

\section{Положительность не выбирает ориентацию}

Существуют два положительных проекторных гамильтониана
\begin{equation}
 h_{\mathbf1\,\min}=P_{\mathbf3}=Q+\frac34I_4,
 \qquad
 h_{\mathbf3\,\min}=P_{\mathbf1}=-Q+\frac14I_4.
 \label{eq:v8-baryon-c0-central-gap-positive-orientation-witnesses}
\end{equation}
Первый оставляет синглет нижним уровнем и имеет $\lambda=+1$, второй
оставляет нижним триплет и имеет $\lambda=-1$. При $\beta>0$ их отклонения
от счётного веса факторизуются как
\begin{equation}
 p_+(\beta)-\frac14
 =\frac{3(e^\beta-1)}{4(e^\beta+3)}>0,
 \qquad
 p_-(\beta)-\frac14
 =-\frac{3(e^\beta-1)}{4(3e^\beta+1)}<0.
 \label{eq:v8-baryon-c0-central-gap-opposite-weight-witnesses}
\end{equation}
Следовательно, самосопряжённость, ограниченность и положительность после
добавления скаляра не выбирают знак щели.

Для произвольного коэффициента равновесный параметр остаётся
\begin{equation}
 \theta=\beta\lambda,
 \qquad \lambda\in\mathbb R.
 \label{eq:v8-baryon-c0-central-gap-free-coefficient}
\end{equation}
Даже разрешение произвольного полинома от семейного оператора Казимира не
снимает свободу. Для
\begin{equation}
 f(\mathcal C_2)=a_0I_4+a_1\mathcal C_2
 +a_2\mathcal C_2^2+a_3\mathcal C_2^3
 \label{eq:v8-baryon-c0-central-gap-casimir-polynomial}
\end{equation}
получаем
\begin{equation}
 \Delta_f=f(2)-f(0)=2a_1+4a_2+8a_3.
 \label{eq:v8-baryon-c0-central-gap-polynomial-coefficient}
\end{equation}
Коэффициенты $a_k$ не фиксируются самим спектром $\{0,2\}$.

Кроме того, чисто квадратичный чётный функционал теряет ориентацию:
\begin{equation}
 (\lambda Q)^2=(-\lambda Q)^2,
 \qquad \Tr Q^2=\frac34.
 \label{eq:v8-baryon-c0-central-gap-quadratic-sign-blindness}
\end{equation}
Для выбора знака потребовался бы линейный или иной нечётный по $Q$ член с
выведенным коэффициентом.

\section{Почему условная форма ещё не принадлежит текущему родителю}

Семейное действие и grading-совместимый endpoint были получены лишь как
минимальные расширения:
\begin{equation}
 N_{SO(3)\,\rm origin}=0/5,
 \qquad
 N_{\rm endpoint\ extension}=0/3.
 \label{eq:v8-baryon-c0-central-gap-unoriginated-structure-ledgers}
\end{equation}
Старое действие не содержит новую triplet-линию, полного $M_3$-угла или
семейного ковариантного кадра. Поэтому операторы $Q$, $\mathcal C_2$ и
$P_{\mathbf3}$ являются согласованными структурами условной архитектуры,
но не ограничениями уже существующего родителя.

Разделим два результата:
\begin{equation}
 N_{\rm conditional\ shape}=3/3,
 \qquad
 N_{\rm physical\ coefficient}=0/5.
 \label{eq:v8-baryon-c0-central-gap-shape-coefficient-ledger}
\end{equation}
Три источника формы --- grading, семейный оператор Казимира и endpoint-
центр --- совпадают modulo скаляры. Проверенные источники физического
коэффициента --- старое действие, линейное отождествление grading с
энергией, нормировка Казимира, положительность и чётный квадратичный
функционал --- не выбирают $\lambda$ и его знак.

Итоговый статус имеет вид
\begin{equation}
 \text{центральное направление: условно единственно},
 \qquad
 \theta=\beta\lambda:\ \text{не выведено}.
 \label{eq:v8-baryon-c0-central-gap-parent-origin-verdict}
\end{equation}

\begin{theorem}[Совпадение grading- и Casimir-направлений щели]
На условном носителе $\mathbf1\oplus\mathbf3$ центрированные grading и
семейный квадратичный оператор Казимира совпадают и равны $2Q$. Поэтому
форма всякого центрального гамильтониана единственна modulo скаляр и
задаётся $Q$. Ни нормировка представления, ни положительность, ни чётный
квадратичный parent не выбирают коэффициент $\lambda$ или его знак.
\end{theorem}

\begin{nogocriterion}[Кинематический grading не является энергией]
Нельзя объявлять $h=P_{\mathbf3}$ родительским гамильтонианом только потому,
что $P_{\mathbf3}=(I+\Gamma_4)/2=\mathcal C_2/2$. Это точное равенство
операторов классифицирует допустимое направление, но не создаёт линейный
энергетический член, его коэффициент, температуру или bath-норму.
\end{nogocriterion}

\begin{protocol}\begin{enumerate}
\item центрированное направление: \textbf{$Q=P_{\mathbf3}-3I_4/4$};
\item семейный оператор Казимира: \textbf{$\mathcal C_2=2P_{\mathbf3}$};
\item центрированные grading и Казимир: \textbf{совпадают};
\item размерность центральной seed-алгебры: \textbf{$2$};
\item условная форма щели: \textbf{единственна modulo скаляр};
\item положительные свидетели двух знаков: \textbf{$P_{\mathbf3}$ и $P_{\mathbf1}$};
\item квадратичный parent видит знак: \textbf{нет};
\item условных источников формы: \textbf{$3/3$};
\item источников физического коэффициента: \textbf{$0/5$};
\item $\theta$ выведено: \textbf{нет};
\item следующий гейт: \textbf{селектор коэффициента центральной щели}.
\end{enumerate}\end{protocol}

Машинный сертификат сохранён в
\path{s2t_v8_baryon_c0_singlet_triplet_central_gap_parent_action_origin_gate_results.json}.